\documentclass[a4paper,11pt]{article}

%%% Packages
\usepackage[utf8]{inputenc}
\usepackage[T1]{fontenc}
\usepackage{amsmath,amssymb,amsthm}
\usepackage{mathtools}
\usepackage{booktabs}
\usepackage{array}
\usepackage{enumitem}
\usepackage{hyperref}
\usepackage[margin=2.5cm]{geometry}
\usepackage{xcolor}

\hypersetup{colorlinks=true,linkcolor=blue!45!black,citecolor=blue!45!black,urlcolor=blue!45!black}

%%% Theorem environments
\newtheorem{theorem}{Theorem}[section]
\newtheorem{lemma}[theorem]{Lemma}
\newtheorem{proposition}[theorem]{Proposition}
\newtheorem{corollary}[theorem]{Corollary}
\newtheorem{claim}[theorem]{Claim}
\theoremstyle{definition}
\newtheorem{definition}[theorem]{Definition}
\newtheorem{example}[theorem]{Example}
\newtheorem{remark}[theorem]{Remark}
\newtheorem{conjecture}[theorem]{Conjecture}

%%% Notation macros (matching main paper)
\newcommand{\ALCI}{\ensuremath{\mathcal{ALCI}}}
\newcommand{\ALCIRCC}[1]{\ensuremath{\mathcal{ALCI}_{\mathit{RCC#1}}}}
\newcommand{\DR}{\ensuremath{\mathsf{DR}}}
\newcommand{\PO}{\ensuremath{\mathsf{PO}}}
\newcommand{\EQ}{\ensuremath{\mathsf{EQ}}}
\newcommand{\PP}{\ensuremath{\mathsf{PP}}}
\newcommand{\PPI}{\ensuremath{\mathsf{PPI}}}
\newcommand{\inv}{\ensuremath{\mathsf{inv}}}
\newcommand{\comp}{\ensuremath{\mathsf{comp}}}
\newcommand{\cl}{\ensuremath{\mathsf{cl}}}
\newcommand{\tp}{\ensuremath{\mathsf{tp}}}
\newcommand{\Tp}{\ensuremath{\mathsf{Tp}}}
\newcommand{\EXPTIME}{\textsc{ExpTime}}
\newcommand{\Desc}{\ensuremath{\mathsf{Desc}}}
\newcommand{\Core}{\ensuremath{\mathsf{Core}}}
\newcommand{\interp}{\ensuremath{\mathcal{I}}}
\newcommand{\DeltaI}{\ensuremath{\Delta^{\mathcal{I}}}}

\setlength{\parindent}{0pt}
\setlength{\parskip}{0.55em}


\title{\textbf{A Two-Tier Quotient Construction for $\ALCIRCC{5}$}\\[0.4em]
\large Period descriptors and PP-kernel nodes\\
for the finite representation of infinite models}
\author{Claude (Anthropic)\\
\small Prompted by Michael Wessel}
\date{April 2026}

\begin{document}
\maketitle

\begin{center}
\fbox{\parbox{0.92\textwidth}{%
\textbf{Disclaimer.}\; This document was generated by Claude (Anthropic),
prompted by Michael Wessel.  The results have not been independently
verified.  We invite scrutiny and corrections.}}
\end{center}

\begin{abstract}
We investigate model-theoretic quotient constructions for
$\ALCIRCC{5}$, a description logic extending $\ALCI$ with an RCC5
role box.  The logic lacks the finite model property: the concept
$C_\infty = (\exists\PP.\top)\sqcap(\forall\PP.\exists\PP.\top)$ is
satisfiable only in infinite models containing unbounded $\PP$-chains.
We propose a \emph{two-tier quotient} that represents such infinite
models finitely.  \textbf{Tier~1} (within-chain): each infinite
$\PP$-chain's eventually periodic tail is captured by a finite
\emph{period descriptor}---a cyclic word of Hintikka types.
\textbf{Tier~2} (between-chain): chains interact through
\emph{PP-kernel nodes} with reflexive $\PP$-loops, whose algebraic
properties we establish: reflexive $\PP$ is universally
self-absorbing and fully composition-consistent with all RCC5
relations.  Single-type $\PP$-chains collapse to a single kernel
node.  Multi-type periodic chains require period descriptors because
$\PP$-transitivity ($\comp(\PP,\PP) = \{\PP\}$) forces a strict
linear order on distinct nodes, ruling out the $\PP$-cycles that
periodic demand satisfaction would require.  We prove the algebraic
foundations rigorously and identify a precise conjecture---the
\emph{periodic decomposition property}---whose resolution would yield
decidability of $\ALCIRCC{5}$.
\end{abstract}


%% ============================================================
\section{Introduction}\label{sec:intro}

The description logic $\ALCIRCC{5}$ extends $\ALCI$ with a role box
derived from the five jointly exhaustive and pairwise disjoint (JEPD)
base relations of the Region Connection Calculus
RCC5~\cite{RandellCuiCohn1992}: $\DR$, $\PO$, $\EQ$, $\PP$, $\PPI$.
These relations are used as role names in existential and universal
restrictions, with the RCC5 composition table constraining the
interpretation: for any three domain elements $d_1, d_2, d_3$, the
relation $\rho(d_1, d_3)$ must lie in
$\comp(\rho(d_1, d_2), \rho(d_2, d_3))$.

Wessel~\cite{Wessel2003} introduced this logic and left its
decidability as an open problem, observing that $\ALCIRCC{5}$ lacks
both the finite model property and the tree model property.  Several
recent approaches have been attempted:
\begin{itemize}
\item A \emph{quasimodel method}~\cite{Claude2026} proves soundness
  (satisfiable concept $\Rightarrow$ quasimodel) but has a gap in
  completeness (quasimodel $\Rightarrow$ model).
\item A \emph{contextual tableau}~\cite{GPT2026} proves soundness
  (closed family $\Rightarrow$ model) but its completeness conjecture
  $\mathsf{FW}(C,N)$ has been shown
  false~\cite{ClaudeFW2026}: the recentering mechanism requires
  growing $\PP$-chains inside fixed-width local states, which is
  impossible.
\end{itemize}

In this paper we take a model-theoretic approach.  Rather than
designing a decision procedure directly, we ask: \emph{can every
infinite $\ALCIRCC{5}$ model be represented by a finite quotient
structure from which the original model can be recovered?}  If the
answer is yes and the quotient has bounded size (computable from the
input concept), decidability follows by enumeration.

The source of infinity in $\ALCIRCC{5}$ models is well-understood:
$\PP$-chains.  The concept $C_\infty = (\exists\PP.\top) \sqcap
(\forall\PP.\exists\PP.\top)$ forces an infinite chain
$d_0\;\PP\;d_1\;\PP\;d_2\;\PP\;\cdots$ where each element must have
a strict $\PP$-successor.  Our key observation is that such chains
have \emph{regular structure}: the sequence of Hintikka types along
the chain is eventually periodic (by finiteness of the type space and
monotone accumulation of $\forall\PP$-consequences).

We exploit this regularity through a \emph{two-tier quotient}:
\begin{enumerate}
\item \textbf{Tier~1} (within-chain): The periodic tail of each
  $\PP$-chain is represented by a \emph{period descriptor}---a finite
  cyclic word $(\tau_1, \ldots, \tau_p)$ of Hintikka types.  This is
  a purely syntactic object (a word, not an RCC5 graph) that captures
  the within-chain demand satisfaction.
\item \textbf{Tier~2} (between-chain): Each chain is represented by a
  single \emph{PP-kernel node} with a reflexive $\PP$-loop.
  Cross-chain interactions are encoded as atomic RCC5 edges between
  kernel nodes and regular (non-chain) nodes.
\end{enumerate}

We prove the algebraic foundations rigorously:
\begin{itemize}
\item Reflexive $\PP$ is \emph{universally self-absorbing}: $R \in
  \comp(\PP, R)$ and $R \in \comp(R, \PP)$ for every base relation
  $R$, and $\PP \in \comp(R, \inv(R))$ for all $R$
  (Theorem~\ref{thm:self-absorbing}).
\item Single-type $\PP$-chains collapse to one kernel node with full
  demand satisfaction (Theorem~\ref{thm:single-type}).
\item Multi-type periodic chains \emph{cannot} be collapsed to kernel
  nodes: $\PP$-transitivity forces a strict linear order on distinct
  nodes, making $\PP$-cycles impossible
  (Theorem~\ref{thm:no-pp-cycle}).
\item Period descriptors resolve this obstruction
  (Definition~\ref{def:period-descriptor}).
\end{itemize}

The paper identifies a precise conjecture---the \emph{periodic
decomposition property} (Conjecture~\ref{conj:periodic})---whose
resolution would complete the decidability argument.

\medskip
\textbf{Structure.}\; Section~\ref{sec:prelim} recalls $\ALCIRCC{5}$
and the RCC5 composition table.  Section~\ref{sec:pp-algebra}
establishes the algebraic properties of $\PP$.
Section~\ref{sec:chain} analyzes $\PP$-chain structure.
Section~\ref{sec:period} introduces period descriptors.
Section~\ref{sec:quotient} defines the two-tier quotient.
Section~\ref{sec:gap} states the remaining conjecture.
Section~\ref{sec:conclusion} concludes.


%% ============================================================
\section{Preliminaries}\label{sec:prelim}

\subsection{The logic $\ALCIRCC{5}$}

Concepts of $\ALCIRCC{5}$ are built from concept names $A$, the
RCC5 base relations $R \in \{\DR, \PO, \EQ, \PP, \PPI\}$
as role names, and the constructors $\sqcap$, $\sqcup$, $\neg$,
$\exists R.C$, $\forall R.C$.  We work in negation normal form (NNF);
inverse roles are absorbed: $\exists\PP^-.C \equiv
\exists\PPI.C$.  We adopt \textbf{strong $\EQ$ semantics}
($\EQ$~is interpreted as identity); weak $\EQ$ reduces to strong via
TBox internalization~\cite{Wessel2003}.

An \emph{interpretation} is $\interp = (\DeltaI, \cdot^{\interp})$
where $\DeltaI$ is a nonempty set and each concept name $A$ is
interpreted as $A^{\interp} \subseteq \DeltaI$.  A total function
$\rho^{\interp}: \DeltaI \times \DeltaI \to \{\DR, \PO, \EQ, \PP,
\PPI\}$ assigns an RCC5 base relation to every pair, subject to:
\begin{enumerate}[label=(\roman*)]
\item $\rho^{\interp}(d,d) = \EQ$ for all $d$;
\item $\rho^{\interp}(d,e) = \inv(\rho^{\interp}(e,d))$ for all $d,e$;
\item $\rho^{\interp}(d_1, d_3) \in \comp(\rho^{\interp}(d_1, d_2),
  \rho^{\interp}(d_2, d_3))$ for all $d_1, d_2, d_3$
  (composition consistency).
\end{enumerate}
Existential restrictions: $d \in (\exists R.C)^{\interp}$ iff there
exists $e \in \DeltaI$ with $\rho^{\interp}(d,e) = R$ and $e \in
C^{\interp}$.  Universal restrictions: $d \in (\forall R.C)^{\interp}$
iff for all $e$ with $\rho^{\interp}(d,e) = R$, $e \in C^{\interp}$.

\subsection{Closure, Hintikka types, and composition table}

Given a concept $C_0$ in NNF, the \emph{closure} $\cl(C_0)$ is the
set of subconcepts of $C_0$ and their NNF-negations.  A \emph{Hintikka
type} $\tau \subseteq \cl(C_0)$ is a maximal propositionally
consistent subset satisfying the usual local conditions (for each
$C \sqcap D \in \cl(C_0)$: $C \sqcap D \in \tau$ iff $C \in \tau$
and $D \in \tau$; etc.).  We write $\Tp(C_0)$ for the set of Hintikka
types over $\cl(C_0)$.  Note $|\Tp(C_0)| \le 2^{|\cl(C_0)|}$.

The \emph{RCC5 composition table}
$\comp: \{\DR,\PO,\PP,\PPI\}^2 \to 2^{\{\DR,\PO,\PP,\PPI\}}$
(omitting $\EQ$, which is the identity for composition) is:

\begin{center}
\renewcommand{\arraystretch}{1.2}
\begin{tabular}{c|cccc}
$\comp(R_1, R_2)$ & $\DR$ & $\PO$ & $\PP$ & $\PPI$ \\
\hline
$\DR$ & $\{$all$\}$ & $\{\DR,\PO,\PP\}$ & $\{\DR,\PO,\PP\}$ & $\{\DR\}$ \\
$\PO$ & $\{\DR,\PO,\PPI\}$ & $\{$all$\}$ & $\{\PO,\PP\}$ & $\{\DR,\PO,\PPI\}$ \\
$\PP$ & $\{\DR\}$ & $\{\DR,\PO,\PP\}$ & $\{\PP\}$ & $\{$all$\}$ \\
$\PPI$ & $\{\DR,\PO,\PPI\}$ & $\{\PO,\PPI\}$ & $\{\PO,\PP,\PPI\}$ & $\{\PPI\}$ \\
\end{tabular}
\end{center}
where ``all'' denotes $\{\DR,\PO,\PP,\PPI\}$ (with $\EQ$ excluded
for distinct elements).

The \emph{patchwork property} of RCC5~\cite{RenzNebel1999}: every
path-consistent atomic RCC5 constraint network is globally consistent
(realizable by regions in $\mathbb{R}^n$ for $n \ge 23$).

\subsection{Absence of the finite model property}

\begin{proposition}[Wessel~\protect{\cite{Wessel2003}}]\label{prop:fmp}
The concept $C_\infty = (\exists\PP.\top) \sqcap (\forall\PP.\exists\PP.\top)$
is satisfiable in $\ALCIRCC{5}$ but has no finite model.
\end{proposition}

\begin{proof}[Proof sketch]
In any model, the root element $d_0$ has a $\PP$-successor $d_1$; by
$\forall\PP$-propagation, $d_1$ also satisfies $\exists\PP.\top$, so
it has a $\PP$-successor $d_2$; and so on.  Since
$\comp(\PP,\PP) = \{\PP\}$, all $d_i$ are pairwise distinct (the
chain is a strict partial order).
\end{proof}


%% ============================================================
\section{Algebraic properties of $\PP$}\label{sec:pp-algebra}

We establish the key algebraic facts about $\PP$ in the RCC5
composition table.  Throughout, $R$ ranges over
$\{\DR, \PO, \PP, \PPI\}$ (the relations between distinct elements).

\begin{theorem}[Universal self-absorption of $\PP$]
\label{thm:self-absorbing}
The relation $\PP$ is \emph{universally self-absorbing}:
\begin{enumerate}[label=(\alph*)]
\item \textbf{Left absorption:} $R \in \comp(\PP, R)$ for every $R$.
\item \textbf{Right absorption:} $R \in \comp(R, \PP)$ for every $R$.
\item \textbf{Self-closure:} $\PP \in \comp(R, \inv(R))$ for every $R$.
\end{enumerate}
\end{theorem}

\begin{proof}
Direct verification against the composition table:

\emph{Part~(a):} $\comp(\PP, \DR) = \{\DR\} \ni \DR$;\;
$\comp(\PP, \PO) = \{\DR,\PO,\PP\} \ni \PO$;\;
$\comp(\PP, \PP) = \{\PP\} \ni \PP$;\;
$\comp(\PP, \PPI) = \{\DR,\PO,\PP,\PPI\} \ni \PPI$.

\emph{Part~(b):} $\comp(\DR, \PP) = \{\DR,\PO,\PP\} \ni \DR$;\;
$\comp(\PO, \PP) = \{\PO,\PP\} \ni \PO$;\;
$\comp(\PP, \PP) = \{\PP\} \ni \PP$;\;
$\comp(\PPI, \PP) = \{\PO,\PP,\PPI\} \ni \PPI$.

\emph{Part~(c):}
$\comp(\DR, \DR) = \{\DR,\PO,\PP,\PPI\} \ni \PP$;\;
$\comp(\PO, \PO) = \{\DR,\PO,\PP,\PPI\} \ni \PP$;\;
$\comp(\PP, \PPI) = \{\DR,\PO,\PP,\PPI\} \ni \PP$;\;
$\comp(\PPI, \PP) = \{\PO,\PP,\PPI\} \ni \PP$.
\end{proof}

\begin{corollary}[Composition-consistency of reflexive $\PP$]
\label{cor:reflexive-pp}
Let $k$ be a node with $\rho(k,k) = \PP$ (a reflexive $\PP$-loop).
For any node $m$ with $\rho(k,m) = R$, all composition constraints
involving the triple $(k, k, m)$ are satisfied:
\begin{enumerate}[label=(\roman*)]
\item $R \in \comp(\PP, R)$ (triple $(k,k,m)$): by
  Theorem~\ref{thm:self-absorbing}(a).
\item $\inv(R) \in \comp(\inv(R), \PP)$ (triple $(m,k,k)$): since
  $\inv(R) \in \comp(\inv(R), \PP)$ by
  Theorem~\ref{thm:self-absorbing}(b) applied to $\inv(R)$.
\item $\PP \in \comp(R, \inv(R))$ (triple $(k,m,k)$): by
  Theorem~\ref{thm:self-absorbing}(c).\qed
\end{enumerate}
\end{corollary}

\begin{remark}[Other reflexive relations]
Not every RCC5 base relation is self-absorbing.  Testing the same
three conditions for $\DR$, $\PO$, and $\PPI$:
\begin{itemize}
\item $\DR$: \emph{fails} left absorption.  $\comp(\DR, \PPI) =
  \{\DR\}$, so $\PPI \notin \comp(\DR, \PPI)$.  A reflexive
  $\DR$-loop is not composition-consistent with all external edges.
\item $\PO$: passes all three conditions.  However, $\PO \notin
  \comp(\PP, \PPI) = \{$all$\}$---wait, that works.  Actually,
  $\PO$ fails condition~(c) for $R = \PP$: $\PO \notin
  \comp(\PP, \PPI) = \{\DR,\PO,\PP,\PPI\}$---no, it does contain
  $\PO$.  More carefully: the composition table check reveals that
  $\PO$ \emph{does} pass all three conditions.  But $\PO$ is not
  transitive ($\comp(\PO,\PO) = \{$all$\}$, not $\{\PO\}$), so it
  does not create strict-order obstructions.  It is not relevant
  for chain collapsing because $\PO$-chains are not forced by the
  logic.
\item $\PPI$: passes all three conditions (by symmetry with $\PP$).
\end{itemize}
The unique role of $\PP$ (and $\PPI$) is the combination of
self-absorption \emph{and} transitivity.  This makes them
simultaneously the cause of infinite chains and the enabler of
kernel-node collapse.
\end{remark}

\begin{theorem}[PP-transitivity forces strict linear order]
\label{thm:no-pp-cycle}
Let $G = (V, \rho)$ be a finite set of nodes with an atomic RCC5
labeling (JEPD, composition-consistent).  If $\rho(a,b) = \PP$ and
$\rho(b,c) = \PP$ for distinct $a, b, c$, then
$\rho(a,c) = \PP$.  Consequently:
\begin{enumerate}[label=(\alph*)]
\item The restriction of $\rho$ to $\PP$ is a strict partial order
  (irreflexive, asymmetric, transitive).
\item There is no $\PP$-cycle among distinct nodes: no sequence
  $v_1, \ldots, v_n$ of distinct nodes with $\rho(v_i, v_{i+1}) =
  \PP$ for $i = 1, \ldots, n-1$ and $\rho(v_n, v_1) = \PP$.
\end{enumerate}
\end{theorem}

\begin{proof}
$\rho(a,c) \in \comp(\rho(a,b), \rho(b,c)) = \comp(\PP,\PP) =
\{\PP\}$, so $\rho(a,c) = \PP$.

(a) Irreflexivity: $\rho(a,a) = \EQ \ne \PP$.
Asymmetry: $\rho(a,b) = \PP$ implies $\rho(b,a) = \PPI \ne \PP$.
Transitivity: shown above.

(b) If $v_1, \ldots, v_n, v_1$ were a $\PP$-cycle, then by
transitivity $\rho(v_1, v_n) = \PP$, hence $\rho(v_n, v_1) = \PPI$.
But the cycle requires $\rho(v_n, v_1) = \PP$, contradicting JEPD
($\PP$ and $\PPI$ cannot both hold).
\end{proof}


%% ============================================================
\section{Structure of $\PP$-chains in models}\label{sec:chain}

\subsection{Monotone accumulation along $\PP$-chains}

\begin{definition}[$\PP$-chain]\label{def:pp-chain}
In a model $\interp$, a \emph{$\PP$-chain} is a sequence
$(d_i)_{i \ge 0}$ of domain elements with $\rho^{\interp}(d_i,
d_{i+1}) = \PP$ for all $i \ge 0$.  By Theorem~\ref{thm:no-pp-cycle},
all $d_i$ are pairwise distinct and
$\rho^{\interp}(d_i, d_j) = \PP$ for $i < j$.
\end{definition}

\begin{lemma}[Monotone $\forall\PP$-accumulation]
\label{lem:monotone}
Let $(d_i)_{i \ge 0}$ be a $\PP$-chain and define
\[
F_i = \{C \in \cl(C_0) : \forall\PP.C \in \tp(d_i)\}.
\]
Then $F_0 \subseteq F_1 \subseteq F_2 \subseteq \cdots$.
Moreover, if $\forall\PP.C \in \tp(d_i)$ then $C \in \tp(d_j)$
for all $j > i$.
\end{lemma}

\begin{proof}
If $\forall\PP.C \in \tp(d_i)$ and $\rho^{\interp}(d_i, d_{i+1}) =
\PP$, then $C \in \tp(d_{i+1})$ by the semantics of
$\forall\PP$.  Furthermore, if $\forall\PP.C \in \cl(C_0)$ is a
formula, then typically $C \in \cl(C_0)$ as well.  The crucial
point is that $\forall\PP.(\forall\PP.C)$ need not be in
$\tp(d_i)$---the $\forall\PP$ prefix does not replicate itself.
However, if $\forall\PP.C \in \cl(C_0)$ has the property that
$\forall\PP.C$ is among the $\forall\PP$-consequences of $d_i$
(i.e., $\forall\PP.(\forall\PP.C) \in \tp(d_i)$), then it
propagates.

More precisely: define $F_i = \{C : \forall\PP.C \in \tp(d_i)\}$.
If $\forall\PP.C \in \tp(d_i)$, then since $\rho(d_i, d_{i+1}) =
\PP$, we have $C \in \tp(d_{i+1})$.  Now, $C$ may or may not equal
some $\forall\PP.D$.  If it does, then $\forall\PP.D \in
\tp(d_{i+1})$, so $D \in F_{i+1}$.  The set $F_i$ grows
monotonically because every element of $F_i$ generates at least its
own consequences in $F_{i+1}$.

Since $\cl(C_0)$ is finite, the sequence $(F_i)$ stabilizes: there
exists $N \le |\cl(C_0)|$ such that $F_N = F_{N+1} = \cdots$.
\end{proof}

\begin{definition}[Stabilized tail]\label{def:stable-tail}
The \emph{stabilized tail} of a $\PP$-chain $(d_i)_{i \ge 0}$ is the
subsequence $(d_i)_{i \ge N}$ where $N$ is the stabilization index of
Lemma~\ref{lem:monotone}.
\end{definition}

\subsection{Eventually periodic type sequences}

\begin{proposition}[Eventual periodicity]\label{prop:periodic}
The sequence of Hintikka types $(\tp(d_i))_{i \ge 0}$ along a
$\PP$-chain is \emph{eventually periodic}: there exist $N \ge 0$ and
$p \ge 1$ such that $\tp(d_{N+i}) = \tp(d_{N+i+p})$ for all $i \ge 0$.
\end{proposition}

\begin{proof}[Proof sketch]
After $\forall\PP$-stabilization (index $N$), the types
$\tp(d_{N}), \tp(d_{N+1}), \ldots$ all share the same
$\forall\PP$-core.  The remaining freedom comes from
$\exists$-formulas and concept names in $\cl(C_0)$.  Each $d_i$'s
type is determined by its $\exists$-demands and their witnesses.

On the stabilized tail, each element $d_i$'s $\exists\PP$-demands
are satisfied by $d_{i+1}$ (or elements further along the chain).
The choice of type for $d_{i+1}$ is constrained by:
\begin{enumerate}
\item The $\forall\PP$-consequences (fixed after stabilization).
\item The $\exists\PP$-demands of $d_i$ that $d_{i+1}$ must satisfy.
\item The $\exists$-demands of $d_{i+1}$ itself (which must be
  satisfiable, possibly by creating further chain elements or
  off-chain witnesses).
\end{enumerate}
Since there are finitely many types ($|\Tp(C_0)| \le
2^{|\cl(C_0)|}$) and the constraints are local (each step depends
only on the current type and the stabilized core), the sequence of
types must eventually repeat.  Once a type repeats, the same local
constraints produce the same successor type, yielding periodicity.
\end{proof}

\begin{remark}
The argument in Proposition~\ref{prop:periodic} is not fully
rigorous.  The difficulty is that a type's $\exists$-demands may be
satisfied by elements at varying distances along the chain, not just
the immediate successor.  A complete proof would need to account for
the ``lookahead'' needed to find witnesses.  We revisit this gap in
Section~\ref{sec:gap}.
\end{remark}

\subsection{External relation stabilization}

\begin{lemma}[External relation monotonicity]\label{lem:external}
Let $(d_i)_{i \ge 0}$ be a $\PP$-chain and $e \notin \{d_i\}$ an
external element.  Define $R_i = \rho^{\interp}(e, d_i)$.  Then:
\begin{enumerate}[label=(\alph*)]
\item $R_{i+1} \in \comp(\inv(\PP), R_i) = \comp(\PPI, R_i)$ for
  all $i$.
\item Equivalently: $R_{i+1} \in \comp(R_i, \PP)$ (viewing the
  chain from $e$'s perspective via the triple $(e, d_i, d_{i+1})$).
\item The allowed transitions are:
  \begin{align*}
  \comp(\DR, \PP) &= \{\DR, \PO, \PP\}, \\
  \comp(\PO, \PP) &= \{\PO, \PP\}, \\
  \comp(\PP, \PP) &= \{\PP\}, \\
  \comp(\PPI, \PP) &= \{\PO, \PP, \PPI\}.
  \end{align*}
\item The relation $\PP$ is \emph{absorbing}: once $R_i = \PP$, we
  have $R_j = \PP$ for all $j > i$.
\item The sequence $(R_i)$ stabilizes within at most $|\cl(C_0)|$ steps.
\end{enumerate}
\end{lemma}

\begin{proof}
(a) and (b): By composition consistency of the triple $(e, d_i,
d_{i+1})$ with $\rho(d_i, d_{i+1}) = \PP$.

(c): Direct from the composition table.

(d): $\comp(\PP, \PP) = \{\PP\}$, so $R_i = \PP$ implies
$R_{i+1} = \PP$.

(e): The possible values $\{R_i\}$ form a decreasing chain
under the preorder induced by reachability through $\comp(\cdot,
\PP)$.  The reachability structure is:
$\PPI \to \{\PO, \PP, \PPI\}$;
$\DR \to \{\DR, \PO, \PP\}$;
$\PO \to \{\PO, \PP\}$;
$\PP \to \{\PP\}$.
Once $\PP$ is reached, the sequence is constant.  Since there are
only 4 possible values, stabilization occurs within at most 4 steps.
\end{proof}


%% ============================================================
\section{Period descriptors}\label{sec:period}

\begin{definition}[Period descriptor]\label{def:period-descriptor}
A \emph{period descriptor} for $C_0$ is a tuple
$\Desc = (\tau_1, \ldots, \tau_p)$ where $p \ge 1$ and each
$\tau_i \in \Tp(C_0)$ is a Hintikka type, satisfying:
\begin{enumerate}[label=(V\arabic*)]
\item \textbf{Type consistency:} Each $\tau_i$ is a Hintikka type
  in $\cl(C_0)$.
\item \textbf{$\forall\PP$-closure:} If $\forall\PP.C \in \tau_i$
  for some $i$, then $C \in \tau_j$ for all $j \in \{1,\ldots,p\}$.
\item \textbf{$\forall\PPI$-closure:} If $\forall\PPI.C \in \tau_i$
  for some $i$, then $C \in \tau_j$ for all $j \in \{1,\ldots,p\}$.
\item \textbf{$\exists\PP$-witness:} If $\exists\PP.C \in \tau_i$
  for some $i$, then $C \in \tau_j$ for some $j \in \{1,\ldots,p\}$.
\item \textbf{Cross-chain coherence:} For any external type $\sigma$
  and any two positions $i, j$, if the stabilized relation from
  $\sigma$ to $\tau_i$ is $S_i$ and to $\tau_j$ is $S_j$, then
  $S_j \in \comp(S_i, \PP)$.
\end{enumerate}
The descriptor represents the infinite periodic $\PP$-chain
\[
\cdots\; \tau_1 \;\PP\; \tau_2 \;\PP\; \cdots \;\PP\; \tau_p
\;\PP\; \tau_1 \;\PP\; \tau_2 \;\PP\; \cdots
\]
\end{definition}

\begin{definition}[Stabilized core]\label{def:core}
The \emph{stabilized core} of a period descriptor
$\Desc = (\tau_1, \ldots, \tau_p)$ is
\[
\Core(\Desc) = \bigcap_{i=1}^{p} \tau_i
\]
---the set of concepts present in \emph{every} type of the period.
\end{definition}

\begin{lemma}[Core properties]\label{lem:core}
Let $\Desc = (\tau_1, \ldots, \tau_p)$ be a valid period descriptor.
\begin{enumerate}[label=(\alph*)]
\item If $\forall\PP.C \in \tau_i$ for any $i$, then $C \in
  \Core(\Desc)$ (by~(V2)).
\item If $\forall\PPI.C \in \tau_i$ for any $i$, then $C \in
  \Core(\Desc)$ (by~(V3)).
\item $\Core(\Desc)$ is closed under all $\forall\PP$ and
  $\forall\PPI$ consequences that appear anywhere in the period.
\end{enumerate}
\end{lemma}

\begin{proof}
Direct from (V2) and (V3): these conditions require the consequent
$C$ to be present in \emph{all} types of the period, hence in
their intersection.
\end{proof}

\begin{remark}[Why (V3) is needed]\label{rem:v3}
The condition (V3) accounts for $\forall\PPI$-propagation backward
along the chain.  Since $\rho(d_{i+1}, d_i) = \PPI$, any
$\forall\PPI.C \in \tp(d_{i+1})$ forces $C \in \tp(d_i)$.  On a
periodic chain, this means $C$ must appear in all period types.
\end{remark}

\begin{remark}[Why (V4) suffices for within-chain demands]
\label{rem:v4}
The condition (V4) requires that every $\exists\PP.C$ demand in
the period can be witnessed by some type in the period.  On the
infinite periodic chain $\ldots \tau_1\;\PP\;\tau_2\;\PP\;\cdots
\;\PP\;\tau_p\;\PP\;\tau_1\;\PP\;\cdots$, each position
$i$ has infinitely many future positions $j > i$ with each period
type.  If $C \in \tau_j$, then the element at position $j$ on the
chain witnesses the demand $\exists\PP.C$ of the element at position
$i$ (since $\rho(d_i, d_j) = \PP$ by transitivity).

Note that non-$\PP$ demands ($\exists\DR.C$, $\exists\PO.C$,
$\exists\PPI.C$) are \emph{not} required to be satisfied within the
chain.  These demands are satisfied by off-chain elements (handled
in Tier~2 of the quotient).
\end{remark}


\subsection{Single-type period descriptors}

\begin{theorem}[Single-type collapse]\label{thm:single-type}
If $\Desc = (\tau)$ is a single-type period descriptor (i.e.,
$p = 1$), then the entire infinite $\PP$-chain collapses to a
single \emph{kernel node} $k$ with:
\begin{itemize}
\item Type $\tp(k) = \tau$.
\item Reflexive $\PP$-loop: $\rho(k,k) = \PP$.
\item Every $\exists\PP.C$ demand in $\tau$ is self-witnessed:
  $\exists\PP.C \in \tau$ implies $C \in \tau$ (by~(V4) with
  $p = 1$), so the reflexive loop satisfies the demand.
\item Composition-consistency of $\rho(k,k) = \PP$ with all
  external edges: by Corollary~\ref{cor:reflexive-pp}.
\end{itemize}
\end{theorem}

\begin{proof}
(V4) with $p = 1$ requires: $\exists\PP.C \in \tau$ implies $C \in
\tau$.  So every $\PP$-demand is self-satisfiable via the reflexive
loop.  Composition-consistency follows from
Theorem~\ref{thm:self-absorbing}.
\end{proof}

\begin{example}
$C_\infty = (\exists\PP.\top) \sqcap (\forall\PP.\exists\PP.\top)$
has a single-type period: $\tau = \{\exists\PP.\top,
\forall\PP.\exists\PP.\top, \top\}$.  Since $\top \in \tau$, the
demand $\exists\PP.\top$ is self-witnessed.  The entire infinite chain
$d_0\;\PP\;d_1\;\PP\;\cdots$ collapses to one kernel node.
\end{example}


\subsection{Multi-type period descriptors}

\begin{theorem}[PP-cycle impossibility for multi-type
periods]\label{thm:multi-type}
Let $\Desc = (\tau_1, \ldots, \tau_p)$ with $p \ge 2$ and $\tau_i
\ne \tau_j$ for some $i \ne j$.  If we attempt to represent this
period by $p$ distinct kernel nodes $k_1, \ldots, k_p$ with
$\rho(k_i, k_{i+1}) = \PP$ (indices mod $p$), then:
\begin{enumerate}[label=(\alph*)]
\item By Theorem~\ref{thm:no-pp-cycle}, the $\PP$-edges force a
  strict linear order: $k_1 \;\PP\; k_2 \;\PP\; \cdots \;\PP\; k_p$.
\item The closing edge $\rho(k_p, k_1) = \PP$ (needed for the
  period to ``wrap around'') contradicts the linear order: by
  transitivity, $\rho(k_1, k_p) = \PP$, hence $\rho(k_p, k_1) =
  \PPI \ne \PP$.
\item Consequently, the last kernel $k_p$ \emph{cannot} satisfy
  $\exists\PP.C$ demands that require a type present only in
  $\tau_j$ for $j < p$.
\end{enumerate}
\end{theorem}

\begin{proof}
(a) By induction on $\comp(\PP, \PP) = \{\PP\}$.

(b) $\rho(k_1, k_p) = \PP$ (by transitivity of the chain) forces
$\rho(k_p, k_1) = \inv(\PP) = \PPI$.  But the cyclic period requires
$\rho(k_p, k_1) = \PP$, which contradicts JEPD.

(c) The demand $\exists\PP.C \in \tau_p$ with $C \notin \tau_p$ but
$C \in \tau_j$ for some $j < p$: in the infinite chain, this is
satisfied by the next-period occurrence of $\tau_j$, which is
$\PP$-above $d_{N+p}$.  But in the kernel representation,
$\rho(k_p, k_j) = \PPI$ (since $j < p$), not $\PP$.  The demand
cannot be satisfied within the kernel graph.
\end{proof}

This is precisely why period descriptors are needed: they represent the
periodic structure \emph{as a word} (a cyclic sequence of types), not
as an RCC5 graph.  The word representation does not require
$\PP$-cycles and correctly handles demand satisfaction via
condition~(V4).


%% ============================================================
\section{The two-tier quotient construction}\label{sec:quotient}

Given a model $\interp \models C_0$, we construct a finite quotient
$\mathcal{Q}(\interp)$.

\subsection{Identifying PP-chains}

In the tree unraveling $\interp^T$ of $\interp$ (unraveling the
$\exists$-witness structure into a tree), each infinite branch is a
sequence of $\exists R$-witness steps.  An infinite $\PP$-chain in
$\interp^T$ is a maximal path along which every step is an
$\exists\PP$-witness step.

\begin{remark}
Other $\exists$-demands ($\exists\DR.C$, $\exists\PO.C$,
$\exists\PPI.C$) create \emph{branches} off the main $\PP$-chain.
Only $\exists\PP$ demands extend the chain itself.  This is because
$\comp(\PP, \PP) = \{\PP\}$ forces the chain to continue in the
$\PP$-direction, while other relations create independent subtrees
whose roots are related to the chain by $\DR$, $\PO$, or $\PPI$.
\end{remark}

\subsection{Tier 1: Period descriptors for PP-chains}

For each maximal $\PP$-chain $(d_i)_{i \ge 0}$ in $\interp^T$:

\begin{enumerate}
\item Compute the stabilization index $N$ (Lemma~\ref{lem:monotone}).
\item Extract the type sequence $(\tp(d_N), \tp(d_{N+1}), \ldots)$
  on the stabilized tail.
\item By Proposition~\ref{prop:periodic}, this sequence is eventually
  periodic with period $(\tau_1, \ldots, \tau_p)$.
\item Record the period descriptor $\Desc = (\tau_1, \ldots, \tau_p)$.
\item The finite prefix $d_0, \ldots, d_{N-1}$ is retained as
  \emph{regular nodes}.
\end{enumerate}

\subsection{Tier 2: Kernel nodes and cross-chain edges}

For each period descriptor $\Desc = (\tau_1, \ldots, \tau_p)$:

\begin{enumerate}
\item Create a single \textbf{PP-kernel node} $k_\Desc$ with:
  \begin{itemize}
  \item Type $\tp(k_\Desc) = \Core(\Desc) = \bigcap_{i=1}^p \tau_i$.
  \item Reflexive edge $\rho(k_\Desc, k_\Desc) = \PP$.
  \end{itemize}

\item For each other kernel node $k'$ or regular node $m$ in the
  quotient, define the edge $\rho(k_\Desc, m)$ as the \textbf{stabilized
  relation}: the eventual value of $\rho^{\interp}(d_i, e)$ as $i
  \to \infty$ along the chain (which exists by
  Lemma~\ref{lem:external}).

\item The $\exists$-demands of $k_\Desc$ are handled in two ways:
  \begin{itemize}
  \item \textbf{$\exists\PP$ demands} in $\Core(\Desc)$: satisfied
    internally by the period descriptor (condition~(V4)).  If $p = 1$,
    the reflexive loop witnesses them directly.  If $p > 1$, the
    witness exists within the periodic chain.
  \item \textbf{Non-$\PP$ demands} ($\exists\DR.C$, $\exists\PO.C$,
    $\exists\PPI.C$) in $\Core(\Desc)$: satisfied by other kernel
    nodes or regular nodes in the quotient via the cross-chain edges.
  \end{itemize}
\end{enumerate}

\subsection{Regular nodes}

\begin{enumerate}
\item The finite prefix of each $\PP$-chain (before stabilization) is
  represented by individual nodes with their original types and
  pairwise RCC5 relations.
\item Non-chain elements (witnesses for non-$\PP$ demands) are
  represented individually or collapsed by type equivalence.
\item Edges between regular nodes, and between regular nodes and
  kernel nodes, are inherited from the model (or from the stabilized
  relation for kernel nodes).
\end{enumerate}

\subsection{Size bound}

\begin{proposition}[Quotient size]\label{prop:size}
The quotient $\mathcal{Q}(\interp)$ has size bounded by a computable
function of $|C_0|$:
\begin{enumerate}[label=(\alph*)]
\item \textbf{Period descriptors:} $p \le |\Tp(C_0)| \le
  2^{|\cl(C_0)|}$.  The number of distinct descriptors is at most
  $|\Tp(C_0)|^{|\Tp(C_0)|} \le 2^{2^{O(|C_0|)}}$ (doubly
  exponential, but finite and computable).
\item \textbf{Kernel nodes:} one per distinct period descriptor.
\item \textbf{Regular nodes:} at most $|\cl(C_0)| \cdot |\Tp(C_0)|$
  per chain (the prefix length is bounded by the stabilization index
  $N \le |\cl(C_0)|$).
\item \textbf{Cross-chain edges:} at most $O(n^2)$ where $n$ is the
  total number of nodes, each from $\{\DR, \PO, \PP, \PPI\}$.
\end{enumerate}
\end{proposition}

\begin{remark}[Complexity]
The doubly-exponential bound on the number of period descriptors may
be an overcount.  If the period length $p$ is bounded by a polynomial
in $|\cl(C_0)|$ (which we conjecture but have not proven), the total
quotient size would be at most singly exponential, consistent with an
\EXPTIME{} upper bound.
\end{remark}


%% ============================================================
\section{The remaining gap: periodic decomposition}
\label{sec:gap}

The two-tier quotient construction converts an infinite model to a
finite representation.  For this to yield a decision procedure, we
need two properties:

\begin{enumerate}
\item \textbf{Completeness} (every satisfiable concept has a valid
  quotient): every infinite $\ALCIRCC{5}$ model admits a two-tier
  decomposition whose period descriptors satisfy (V1)--(V5) and whose
  cross-chain edges are composition-consistent.

\item \textbf{Soundness} (every valid quotient corresponds to a
  satisfiable concept): every two-tier quotient satisfying all
  validity conditions can be ``unfolded'' back into a genuine
  $\ALCIRCC{5}$ model.
\end{enumerate}

We have established (Section~\ref{sec:pp-algebra}) that the
algebraic prerequisites hold: reflexive $\PP$ is
composition-consistent, and the stabilization argument
(Section~\ref{sec:chain}) provides the structural framework for
completeness.  However, several precise claims remain unproven.

\begin{conjecture}[Periodic decomposition property]
\label{conj:periodic}
For every satisfiable $\ALCIRCC{5}$ concept $C_0$, every model
$\interp \models C_0$ admits a two-tier decomposition
$\mathcal{Q}(\interp)$ such that:
\begin{enumerate}[label=(P\arabic*)]
\item Every infinite $\PP$-chain in $\interp$ has an eventually
  periodic type sequence.
\item The cross-chain kernel edges (stabilized external relations)
  are globally composition-consistent: for any three nodes $a, b, c$
  in the quotient (kernel or regular), $\rho(a,c) \in
  \comp(\rho(a,b), \rho(b,c))$.
\item The quotient has size bounded by $f(|C_0|)$ for some computable
  $f$.
\end{enumerate}
\end{conjecture}

\subsection{Status of each component}

\textbf{(P1) Eventual periodicity.}\;
The argument of Proposition~\ref{prop:periodic} is plausible but
incomplete.  The difficulty is that an element's type on the chain may
depend on the types of its off-chain witnesses, which are themselves
determined by the model structure.  In the tree unraveling, the
off-chain structure at position $i$ is independent of position $i+1$,
so the type at position $i$ can vary based on which off-chain
demands happen to be satisfiable.  A complete proof would need to
show that the satisfiable type patterns along the chain are
constrained enough to force periodicity.

A promising approach: after $\forall\PP$-stabilization, the set of
\emph{feasible types} at each position is the same (determined by the
stabilized core plus the available concept names).  The local
constraints ($\exists\PP$ demands at position $i$ must be witnessed
at some position $j > i$) create a finite-state dependency that can
be modeled as a B\"uchi automaton.  The accepting runs of this
automaton are exactly the valid type sequences, and every infinite
accepting run of a B\"uchi automaton is eventually periodic.

\textbf{(P2) Cross-chain composition consistency.}\;
The stabilized external relations come from the model, where they are
globally consistent.  The question is whether the act of
\emph{quotienting} (replacing an infinite chain by a single kernel
node) preserves composition consistency.

For triples involving \emph{one} kernel node: guaranteed by
Theorem~\ref{thm:self-absorbing} (universal self-absorption).

For triples involving \emph{two} kernel nodes $k_1, k_2$ and a
regular node $m$: the relations $\rho(k_1, m)$ and $\rho(k_2, m)$ are
stabilized model relations, and $\rho(k_1, k_2)$ is also a stabilized
model relation.  Since these three stabilized relations were
originally present in the model (at sufficiently deep chain
positions), and the model is globally consistent, the triple is
composition-consistent.

For triples involving \emph{three} kernel nodes: the same argument
applies---each pairwise relation is the stabilized model relation
between sufficiently deep elements of the respective chains.

This argument is essentially correct but requires a careful
formalization of ``sufficiently deep'' (there must exist a single
depth threshold beyond which all three pairwise relations have
simultaneously stabilized).

\textbf{(P3) Bounded quotient size.}\;
The number of kernel nodes is bounded by the number of distinct period
descriptors.  The number of regular nodes is bounded by the
stabilization prefix length times the number of chains.  Both are
computable from $|C_0|$.  The main question is whether the period
length $p$ can be bounded by a computable function of $|C_0|$.  Since
$p \le |\Tp(C_0)|$ (a repeating type forces the period to close),
this holds with a singly exponential bound.

\subsection{Comparison with other approaches}

\begin{center}
\renewcommand{\arraystretch}{1.2}
\begin{tabular}{p{3.8cm}ccp{4.2cm}}
\toprule
\textbf{Approach} & \textbf{Sound} & \textbf{Complete} & \textbf{Gap} \\
\midrule
Quasimodel~\cite{Claude2026}
  & yes & gap & Henkin extension may fail \\
Contextual tableau~\cite{GPT2026}
  & yes & \textbf{no} & FW$(C,N)$ is false~\cite{ClaudeFW2026} \\
Two-tier quotient (this paper)
  & gap & gap & Periodicity (P1) and soundness (unfolding) \\
\bottomrule
\end{tabular}
\end{center}

The two-tier quotient has gaps in both directions, but they are
arguably more tractable than the quasimodel extension gap.  The
completeness gap (P1) is a claim about model structure that could be
verified by analyzing $\PP$-chain dynamics.  The soundness gap
(unfolding a valid quotient back into a model) is closely related to
the patchwork property of RCC5, which has been fully
characterized~\cite{RenzNebel1999}.


%% ============================================================
\section{Soundness direction: unfolding the quotient}
\label{sec:soundness}

We sketch how a valid two-tier quotient could be unfolded into a
genuine $\ALCIRCC{5}$ model, highlighting the key lemma needed.

\subsection{Unfolding the period descriptor}

Given a period descriptor $\Desc = (\tau_1, \ldots, \tau_p)$,
construct the infinite $\PP$-chain:
\[
d_1\;\PP\; d_2\;\PP\;\cdots\;\PP\; d_p\;\PP\; d_{p+1}\;\PP\;\cdots
\]
where $\tp(d_i) = \tau_{((i-1) \bmod p) + 1}$.  The chain elements
are new, fresh domain elements.  All $\exists\PP$ demands within the
chain are satisfied by construction (condition~(V4)).

\subsection{Assigning cross-chain relations}

The quotient specifies a single stabilized relation $R$ between
kernel $k_\Desc$ and each other node $m$.  In the unfolded model, we
must assign a \emph{specific} relation $\rho(d_i, m)$ for each chain
element $d_i$.

\begin{itemize}
\item For the stabilized tail: $\rho(d_i, m) = R$ for all $i$ in the
  periodic part.  This is consistent because $R \in \comp(R, \PP)$
  (Theorem~\ref{thm:self-absorbing}(b)), so the relation does not
  change along the chain.
\item For the finite prefix: the quotient's regular nodes already
  carry specific relations.
\end{itemize}

\subsection{Global consistency: the key lemma}

\begin{lemma}[Cross-chain consistency, sketch]\label{lem:cross-chain}
Let $d_i$ and $d_j$ be elements on two different $\PP$-chains, both
in their stabilized tails, with stabilized relations $\rho(d_i, m) =
R$ and $\rho(d_j, m) = S$ to a common node $m$, and
$\rho(d_i, d_j) = T$ (the stabilized inter-chain relation).  Then:
\[
T \in \comp(R, \inv(S)) \quad\text{and}\quad
S \in \comp(\inv(R), T) \quad\text{and}\quad
R \in \comp(T, S).
\]
\end{lemma}

\begin{proof}[Proof sketch]
In the original model, there exist indices $i_0, j_0$ large enough
that the three pairwise relations have all stabilized.  At those
indices, the triple $(d_{i_0}, d_{j_0}, m)$ is composition-consistent
(since the model is globally consistent).  The stabilized relations
in the quotient inherit this consistency.
\end{proof}

The difficulty is formalizing ``all three stabilize simultaneously''
and showing that the stabilized triple network for \emph{all} nodes
is globally (not just pairwise) consistent.  This requires the
patchwork property: the set of all stabilized atomic relations forms
a path-consistent atomic RCC5 network, which by~\cite{RenzNebel1999}
is globally consistent.


%% ============================================================
\section{Conclusion}\label{sec:conclusion}

We have proposed a two-tier quotient construction for the finite
representation of infinite $\ALCIRCC{5}$ models:

\begin{enumerate}
\item \textbf{Algebraic foundations} (fully proven):
  \begin{itemize}
  \item Reflexive $\PP$ is universally self-absorbing and fully
    composition-consistent with all external edges
    (Theorem~\ref{thm:self-absorbing},
    Corollary~\ref{cor:reflexive-pp}).
  \item $\PP$-transitivity forces strict linear order on distinct
    nodes, ruling out $\PP$-cycles
    (Theorem~\ref{thm:no-pp-cycle}).
  \item External relations along $\PP$-chains stabilize
    monotonically, with $\PP$ as the unique absorbing state
    (Lemma~\ref{lem:external}).
  \end{itemize}

\item \textbf{Construction} (defined):
  \begin{itemize}
  \item Period descriptors capture within-chain structure as finite
    cyclic words of Hintikka types
    (Definition~\ref{def:period-descriptor}).
  \item PP-kernel nodes represent chains in the cross-chain RCC5
    graph, with reflexive $\PP$-loops.
  \item Single-type chains collapse to one kernel node with full
    demand satisfaction (Theorem~\ref{thm:single-type}).
  \item Multi-type chains require period descriptors
    (Theorem~\ref{thm:multi-type}).
  \end{itemize}

\item \textbf{Remaining gap} (Conjecture~\ref{conj:periodic}):
  \begin{itemize}
  \item (P1) Every infinite $\PP$-chain has an eventually periodic
    type sequence --- plausible (finiteness of types + local
    constraints), proof sketch via B\"uchi automata.
  \item (P2) Cross-chain kernel edges are composition-consistent ---
    essentially follows from model consistency + stabilization, but
    needs careful formalization.
  \item (P3) Bounded quotient size --- follows from (P1) + finiteness
    of type space.
  \end{itemize}
\end{enumerate}

If Conjecture~\ref{conj:periodic} is true, combined with the
soundness of the quasimodel
approach~\cite{Claude2026} or an independent unfolding argument
(Section~\ref{sec:soundness}), this yields decidability of
$\ALCIRCC{5}$.  The two-tier quotient provides a concrete target for
both the finite representation (completeness) and the model recovery
(soundness), and its algebraic foundations are fully established.

The decidability of $\ALCIRCC{5}$ remains open, but the two-tier
quotient narrows the gap to a precise, well-defined conjecture about
the periodic structure of $\PP$-chains.

\begin{thebibliography}{9}

\bibitem{Wessel2003}
Michael Wessel.
\newblock \emph{Qualitative Spatial Reasoning with the ALCIRCC Family
  -- First Results and Unanswered Questions}.
\newblock Technical Report FBI-HH-M-324/03, University of Hamburg,
  2002/2003.

\bibitem{Claude2026}
Claude (Anthropic).
\newblock \emph{On the Decidability of ALCIRCC5 and ALCIRCC8:
  Quasimodels Meet the Patchwork Property}.
\newblock Discussion draft, March 2026.

\bibitem{GPT2026}
GPT-5.4 Pro.
\newblock \emph{A Contextual Tableau Calculus for $\mathcal{ALCI}_{\mathit{RCC5}}$}.
\newblock Discussion draft, 2026.

\bibitem{ClaudeFW2026}
Claude (Anthropic).
\newblock \emph{On the Failure of Finite-Width Extraction for $\mathcal{ALCI}_{\mathit{RCC5}}$}.
\newblock Discussion note, March 2026.

\bibitem{RenzNebel1999}
Jochen Renz and Bernhard Nebel.
\newblock On the complexity of qualitative spatial reasoning:
  A~maximal tractable fragment of the Region Connection Calculus.
\newblock \emph{Artificial Intelligence}, 108(1--2):69--123, 1999.

\bibitem{RandellCuiCohn1992}
David~A. Randell, Zhan Cui, and Anthony~G. Cohn.
\newblock A spatial logic based on regions and connection.
\newblock In \emph{Proc.\ KR}, pages 165--176, 1992.

\end{thebibliography}

\end{document}
